%% toc: P15 160 160.9
% Source-specific presentation: Feller II, VII.5--VII.6 (pp. 230--234), XV.4 (pp. 512--515),
% XIII.1 (pp. 429--434), XV.2 (pp. 504--505), XV.3 (pp. 507--511), X.8--X.9 (pp. 350--355), XIII.9--XIII.10 (pp. 454--463).
\chapter{Feller's Poisson, Laplace and Analytic Moment Continuations}
\label{sec:p15-w02b-p15-012-feller-continuations}

\Cref{sec:p15-feller-bernstein-hausdorff} follows Feller's finite-difference
route through Bernstein polynomials, Hausdorff's moment theorem and
exchangeability. Feller then continues the same laws-of-large-numbers method
in two further directions and later returns to the moment problem with
characteristic functions. The same transforms reappear in his continuity
theorems and in semigroup theory. This chapter supplies these continuations as
separate presentations, each with its own hypotheses, conclusions and proofs:
\begin{enumerate}
\item the Poisson randomization identity, the generalized Taylor formula and
      its operator form for translations
      \cite[vol.~II, VII.5, pp.~230--232]{Feller1971ProbabilityII};
\item the Poisson inversion formula for Laplace transforms, the inversion of
      Laplace integrals by the gamma kernel, the alternating moment
      inequalities and the moment series for laws on \([0,\infty)\)
      \cite[vol.~II, VII.6, pp.~232--234]{Feller1971ProbabilityII};
\item the characteristic-function proof that a real law whose absolute moments
      grow at most like \((\lambda n)^n\) is determined by its moments
      \cite[vol.~II, XV.4, pp.~512--515]{Feller1971ProbabilityII};
\item the continuity theorems for Laplace transforms and for characteristic
      functions, with Feller's second uniqueness proof, his transform examples and his Fourier, Plancherel and lattice inversion formulas \cite[vol.~II, XIII.1, pp.~429--434; XV.2, pp.~504--505; XV.3, pp.~507--511]{Feller1971ProbabilityII};
\item the exponential formula for continuous semigroups of contractions and
      the abstract version of the Laplace-integral inversion (6.6)
      \cite[vol.~II, X.8--X.9, pp.~350--355; XIII.9--XIII.10, pp.~454--463]{Feller1971ProbabilityII}.
\end{enumerate}
The inputs are a bounded continuous function, a probability law on
\([0,\infty)\) or on the line, and, in the moment results, the moment
sequence alone. The conclusions are limit formulas that recover the function
or the distribution function from the stated data, and uniqueness theorems.
In the continuity theorems the input is a sequence of laws or measures with
their transforms, and in the semigroup results it is a continuous semigroup of
contractions on a space of bounded functions; the conclusions identify the
limit law, or recover the semigroup from its values at small times or from its
Laplace transform. The two moment-determinacy results below are Feller's own routes. They are
distinct from Durrett's analytic propagation proof
(\cref{thm:p15-criteria-analytic-determinacy}) and from the Carleman
criterion (\cref{thm:p15-criteria-carleman}); \cref{thm:p15-w02b-p15-012-xv4-bridge}
gives Feller's own proof of the power-series condition stated in
\cref{thm:p15-feller-lognormal-moments}.

Throughout, \(N_\mu\) denotes a random variable with the Poisson law of mean
\(\mu\geq0\), so \(\mathbb P\{N_\mu=k\}=e^{-\mu}\mu^k/k!\) for
\(k=0,1,\ldots\), with \(N_0=0\). Its mean and variance both equal \(\mu\).

\section{The Poisson identity and the generalized Taylor formula}
\label{sec:p15-w02b-p15-012-taylor}

The normalized difference \(\Delta_h\) and the sequence difference
\(\Delta\) are those of \cref{def:p15-feller-differences}. In particular
\(\Delta_h^r u(x)=h^{-r}\Delta^r a_j\) when \(a_i=u(ih)\) and \(x=jh\).

\begin{hbtheorem}{thm:p15-w02b-p15-012-poisson-identity}{proved=annex:p15-w02b-p15-012-poisson-identity}{Feller's Poisson identity (5.1)--(5.2)}
Let \((a_j)_{j\geq0}\) be a bounded real sequence, \(i\geq0\) an integer and
\(x\) a real number. Both series below converge absolutely and
\[
 \sum_{r=0}^\infty\frac{x^r}{r!}\Delta^r a_i
      = e^{-x}\sum_{j=0}^\infty\frac{x^j}{j!}a_{i+j}.
 \tag{5.1}
\]
Consequently, for a bounded function \(u\) on \([0,\infty)\), \(h>0\) and
\(\theta\geq0\),
\[
 \sum_{r=0}^\infty\frac{\theta^r}{r!}\Delta_h^r u(0)
      = e^{-\theta/h}\sum_{j=0}^\infty u(jh)\frac{(\theta/h)^j}{j!}
      = \mathbb E\,u(hN_{\theta/h}).
 \tag{5.2}
\]
The right side of (5.2) is the expectation of \(u\) under the law of
\(hN_{\theta/h}\), which has mean \(\theta\) and variance \(h\theta\).
\end{hbtheorem}

Feller obtains (5.1) as the Poisson limit of the binomial identity behind the
Bernstein representation, and observes that it also follows by expanding
\(\Delta^r a_i\) and rearranging \cite[vol.~II, VII.5, p.~230 and n.~6]{Feller1971ProbabilityII}.
The annex gives the rearrangement with its absolute-convergence justification.

\begin{hbtheorem}{thm:p15-w02b-p15-012-generalized-taylor}{proved=annex:p15-w02b-p15-012-generalized-taylor}{The generalized Taylor formula (5.3)}
Let \(u\) be bounded and continuous on \([0,\infty)\). For every
\(t\geq0\) and \(\theta>0\),
\[
 \sum_{r=0}^\infty\frac{\theta^r}{r!}\Delta_h^r u(t)\longrightarrow u(t+\theta)
 \qquad(h\to0+).
 \tag{5.3}
\]
For fixed \(t\) the convergence is uniform for \(\theta\) in every bounded
interval \([0,b]\), where the value at \(\theta=0\) is \(u(t)\) for every
\(h\). If \(u\) is uniformly continuous on \([0,\infty)\), the convergence is
uniform in \(t\geq0\) and \(0\leq\theta\leq b\).
If \(u\) is bounded and continuous on the whole line and \(\Delta_h\) is
defined by the same quotient \((u(x+h)-u(x))/h\) for \(h<0\), then (5.3)
also holds for every real \(t\) when \(\theta<0\) and \(h\to0-\).
\end{hbtheorem}

The left side of (5.3) is the Taylor series of \(u\) at \(t\) with derivatives
replaced by difference ratios. The theorem requires no differentiability.
Feller attributes the result to Hille and the proof to M.~Riesz; the
probabilistic reading as a Poisson randomization of the parameter is his
\cite[vol.~II, VII.5, p.~231 and nn.~7--8]{Feller1971ProbabilityII}.

\begin{hbdefinition}{def:p15-w02b-p15-012-translation}{Translation operators and the Poisson interpolation operator}
Let \(C_b\) be the space of bounded continuous real functions on
\([0,\infty)\) with the norm \(\|u\|=\sup_{t\geq0}|u(t)|\). For
\(\theta\geq0\) the translation operator \(T(\theta)\) sends \(u\) to
\(u_\theta\), where \(u_\theta(t)=u(t+\theta)\). Thus \(T(0)=1\) is the
identity operator, \(T(\theta_1)T(\theta_2)=T(\theta_1+\theta_2)\), and
\(\|T(\theta)u\|\leq\|u\|\). For \(h>0\), the normalized difference is the
operator
\[
 \Delta_h=h^{-1}[T(h)-1].
 \tag{5.5}
\]
For \(\theta\geq0\) and \(h>0\), Feller's interpolation operator is
\[
 A_h(\theta)=e^{-\theta h^{-1}}\sum_{k=0}^\infty
               \frac1{k!}\Bigl(\frac\theta h\Bigr)^kT(kh),
 \tag{5.7}
\]
and the operator exponential is the norm-convergent series
\[
 e^{\theta h^{-1}[T(h)-1]}=\sum_{r=0}^\infty\frac{\theta^r}{r!}\Delta_h^r .
\]
Both series converge in operator norm, because
\(\|\Delta_h^r\|\leq(2/h)^r\) and \(\|T(kh)\|\leq1\).
\end{hbdefinition}

\begin{hbtheorem}{thm:p15-w02b-p15-012-exponential-formula}{proved=annex:p15-w02b-p15-012-exponential-formula}{The exponential formula (5.6) for translations}
For \(\theta\geq0\) and \(h>0\), \(e^{\theta h^{-1}[T(h)-1]}=A_h(\theta)\).
For every \(u\in C_b\) and every \(t\geq0\),
\[
 e^{\theta h^{-1}[T(h)-1]}u(t)\longrightarrow T(\theta)u(t)
 \qquad(h\to0+),
 \tag{5.6}
\]
and the convergence holds in the norm of \(C_b\) whenever \(u\) is uniformly
continuous. The convergence is not uniform over the unit ball of \(C_b\):
for each \(\theta>0\) and every \(h>0\) with \(\theta/h\) not an integer,
\[
 \|A_h(\theta)-T(\theta)\|=2 .
\]
\end{hbtheorem}

Formula (5.6) says that the whole family \(T(\theta)\) is recovered from the
behaviour of \(T(h)\) for small \(h\). Feller notes that the same Poisson
interpolation makes the exponential formula plausible for every continuous
semigroup of bounded operators and defers the general proof to X.9
\cite[vol.~II, VII.5, p.~232]{Feller1971ProbabilityII}. The theorem above
is restricted to the translation semigroup that VII.5 itself treats. The
general theorem of X.9 is \cref{thm:p15-w02b-p15-012-semigroup-exponential}.

\begin{hbexample}{def:p15-w02b-p15-012-taylor-examples}{Exact Poisson interpolation of an exponential, and a corner}
For \(u(x)=e^{-x}\), \(\Delta_h u=c_h u\) with \(c_h=(e^{-h}-1)/h\).
Hence, exactly,
\[
 \sum_{r=0}^\infty\frac{\theta^r}{r!}\Delta_h^r u(t)
     =e^{-t}\exp\bigl\{\theta(e^{-h}-1)/h\bigr\},
\]
which differs from \(u(t+\theta)=e^{-t-\theta}\) for every \(h>0\) and
\(\theta>0\), and tends to it as \(h\to0+\) because \(c_h\to-1\).
For \(u(x)=\min(x,1)\), which has no derivative at \(1\), the left side of
(5.3) at \(t=0\) is \(\mathbb E\min(hN_{\theta/h},1)\), and
\cref{thm:p15-w02b-p15-012-generalized-taylor} still gives the limit
\(\min(\theta,1)\). Both computations are verified in the annex under
\cref{thm:p15-w02b-p15-012-generalized-taylor}.
\end{hbexample}

\section{Inversion formulas for Laplace transforms}
\label{sec:p15-w02b-p15-012-laplace}

\begin{hbdefinition}{def:p15-w02b-p15-012-laplace}{Laplace transform of a law on \([0,\infty)\)}
Let \(F\) be a probability distribution concentrated on \([0,\infty)\),
written as a measure \(F\{d\theta\}\) and as the distribution function
\(F(x)=F\{[0,x]\}\). Its Laplace transform is
\[
 \varphi(\lambda)=\int_0^\infty e^{-\lambda\theta}\,F\{d\theta\},
 \qquad\lambda\geq0 .
 \tag{6.2}
\]
Feller defines \(\varphi\) for \(\lambda>0\); the value \(\varphi(0)=1\) is
added here. When they are finite, the moments are
\(\mu_k=\int_0^\infty\theta^k\,F\{d\theta\}\), with \(\mu_0=1\).
For a bounded continuous \(u\) on \([0,\infty)\), the Laplace integral is
\[
 w(\lambda)=\int_0^\infty e^{-\lambda x}u(x)\,dx,\qquad\lambda>0 .
 \tag{6.5}
\]
In this section \(\lambda\) is the transform variable.
\end{hbdefinition}

\begin{hbtheorem}{thm:p15-w02b-p15-012-laplace-derivatives}{proved=annex:p15-w02b-p15-012-laplace-derivatives}{Differentiation of Laplace transforms (6.3)}
For \(\lambda>0\) the transform \(\varphi\) has derivatives of every order and
\[
 (-1)^k\varphi^{(k)}(\lambda)=\int_0^\infty e^{-\lambda\theta}\theta^k\,F\{d\theta\},
 \qquad k=0,1,\ldots .
 \tag{6.3}
\]
Likewise \(w\) has derivatives of every order on \(\lambda>0\), and
\(w^{(k)}(\lambda)=\int_0^\infty(-x)^ke^{-\lambda x}u(x)\,dx\).
\end{hbtheorem}

\begin{hbtheorem}{thm:p15-w02b-p15-012-laplace-inversion}{proved=annex:p15-w02b-p15-012-laplace-inversion}{The Poisson inversion formula (6.1) and (6.4)}
For \(\theta\geq0\) and real \(x\neq\theta\), as \(\lambda\to\infty\),
\[
 e^{-\lambda\theta}\sum_{k\leq\lambda x}\frac{(\lambda\theta)^k}{k!}
   =\mathbb P\{N_{\lambda\theta}\leq\lambda x\}
   \longrightarrow
   \begin{cases}0,&\theta>x,\\ 1,&\theta<x.\end{cases}
 \tag{6.1}
\]
Consequently, at every point \(x\) of continuity of \(F\),
\[
 \sum_{k\leq\lambda x}\frac{(-1)^k}{k!}\lambda^k\varphi^{(k)}(\lambda)
   \longrightarrow F(x)\qquad(\lambda\to\infty).
 \tag{6.4}
\]
In particular, a distribution concentrated on \([0,\infty)\) is uniquely
determined by its Laplace transform on any half-line \(\lambda>\lambda_1\).
\end{hbtheorem}

Feller derives (6.1) from the weak law for the Poisson law and notes that it
is the special case of (5.2) for a two-valued \(u\)
\cite[vol.~II, VII.6, p.~232]{Feller1971ProbabilityII}. The two-valued
function is discontinuous, so the annex proves (6.1) directly by Chebyshev's
inequality rather than as an instance of (5.3).

\begin{hbtheorem}{thm:p15-w02b-p15-012-laplace-integral-inversion}{proved=annex:p15-w02b-p15-012-laplace-integral-inversion}{Inversion of Laplace integrals (6.6)}
Let \(u\) be bounded and continuous on \([0,\infty)\) and let \(w\) be its
Laplace integral (6.5). For integers \(n\geq1\) and \(\theta>0\),
\[
 \frac{(-1)^{n-1}}{(n-1)!}\Bigl(\frac n\theta\Bigr)^n w^{(n-1)}(n/\theta)
   =\frac1{(n-1)!}\int_0^\infty u(x)\Bigl(\frac{nx}\theta\Bigr)^{n-1}
          e^{-nx/\theta}\,\frac{n\,dx}\theta
   \longrightarrow u(\theta)
 \tag{6.6}
\]
as \(n\to\infty\), uniformly for \(\theta\) in every interval \(0<\theta\leq b\).
\end{hbtheorem}

The middle expression is the gamma-kernel expectation of VII.1, example (c),
for the law with shape \(n\) and rate \(n/\theta\); its variance is
\(\theta^2/n\) (\cref{thm:p15-feller-kernel-examples}). Feller states that an
abstract version of (6.6) appears in XIII.9
\cite[vol.~II, VII.6, p.~233]{Feller1971ProbabilityII}; it is
\cref{thm:p15-w02b-p15-012-abstract-inversion}.

\begin{hbtheorem}{thm:p15-w02b-p15-012-laplace-moment-inequalities}{proved=annex:p15-w02b-p15-012-laplace-moment-inequalities}{The alternating inequalities (6.7)--(6.9)}
For every integer \(n\geq1\) and every \(t>0\),
\[
 \sum_{k=0}^{2n-1}\frac{(-1)^kt^k}{k!}<e^{-t}<\sum_{k=0}^{2n}\frac{(-1)^kt^k}{k!}.
 \tag{6.8}
\]
If \(F\) has finite moments \(\mu_1,\ldots,\mu_{2n}\), then for every
\(\lambda\geq0\),
\[
 \sum_{k=0}^{2n-1}\frac{(-1)^k\mu_k\lambda^k}{k!}
   \leq\varphi(\lambda)\leq
 \sum_{k=0}^{2n}\frac{(-1)^k\mu_k\lambda^k}{k!}.
 \tag{6.7}
\]
If all moments are finite and \(\sum_k\mu_k\lambda_0^k/k!<\infty\) for some
\(\lambda_0>0\), then
\[
 \varphi(\lambda)=\sum_{k=0}^\infty\frac{(-1)^k\mu_k\lambda^k}{k!},
 \qquad0\leq\lambda\leq\lambda_0 .
 \tag{6.9}
\]
\end{hbtheorem}

\begin{hbtheorem}{thm:p15-w02b-p15-012-laplace-moment-determinacy}{proved=annex:p15-w02b-p15-012-laplace-moment-determinacy}{Moment determinacy on \([0,\infty)\) through the Laplace transform}
Let \(F\) and \(G\) be probability distributions concentrated on
\([0,\infty)\) with the same finite moments \(\mu_k\), \(k\geq0\), and
suppose \(\sum_k\mu_k\lambda_0^k/k!<\infty\) for some \(\lambda_0>0\).
Then the Laplace transforms of \(F\) and \(G\) agree for every
\(\lambda\geq0\), and \(F=G\).
\end{hbtheorem}

Feller passes from (6.9) to all \(\lambda>0\) by appealing to analytic function
theory \cite[vol.~II, VII.6, p.~234]{Feller1971ProbabilityII}. The annex
replaces that appeal by an explicit real-variable propagation: each transform
equals its Taylor series of radius at least \(\lambda_0\) about every point of
\([0,\infty)\), and agreement spreads from \([0,\lambda_0]\) in steps of fixed
length. The uniqueness of \(F\) then follows from (6.4).

\begin{hbtheorem}{thm:p15-w02b-p15-012-laplace-examples}{proved=annex:p15-w02b-p15-012-laplace-examples}{The exponential law, an atom and the lognormal moments}
\begin{enumerate}
\item For the exponential law \(F(x)=1-e^{-x}\), \(x\geq0\),
      \(\varphi(\lambda)=1/(1+\lambda)\), \(\mu_k=k!\), the series (6.9)
      converges for \(0\leq\lambda<1\) to \(1/(1+\lambda)\), and for
      \(x\geq0\) the left side of (6.4) equals
      \(1-(\lambda/(1+\lambda))^{\lfloor\lambda x\rfloor+1}\), which tends to
      \(1-e^{-x}\).
\item Continuity of \(F\) at \(x\) cannot be dropped in (6.4). For
      \(F=\delta_{x_0}\) with \(x_0>0\), along \(\lambda_m=m/x_0\) the left
      side of (6.4) at \(x=x_0\) equals \(\mathbb P\{N_m\leq m\}\), which
      tends to \(1/2\), while \(F(x_0)=1\).
\item The hypothesis of \cref{thm:p15-w02b-p15-012-laplace-moment-determinacy}
      cannot be dropped. The lognormal laws of
      \cref{thm:p15-feller-lognormal-moments} are concentrated on
      \([0,\infty)\), have the common moments \(\mu_k=e^{k^2/2}\), and are
      distinct; for these moments \(\sum_k\mu_k\lambda^k/k!=\infty\) for
      every \(\lambda>0\).
\end{enumerate}
\end{hbtheorem}

\section{Feller's characteristic-function proof of moment determinacy}
\label{sec:p15-w02b-p15-012-xv4}

In this section \(F\) is an arbitrary distribution on the line and
\(\varphi(\zeta)=\int e^{i\zeta x}F\{dx\}\) is its characteristic function,
following Feller's notation in XV. The letter \(\lambda\) now denotes the
growth constant in (4.15), not a transform variable.

\begin{hbdefinition}{def:p15-w02b-p15-012-moments}{Moments and absolute moments (4.4)}
When they exist,
\[
 m_n=\int_{-\infty}^{+\infty}x^n\,F\{dx\},\qquad
 M_n=\int_{-\infty}^{+\infty}|x|^n\,F\{dx\}.
 \tag{4.4}
\]
Thus \(M_{2k}=m_{2k}\), and \(M_n<\infty\) implies that \(M_j<\infty\) and
\(m_j\) exists for \(0\leq j\leq n\).
\end{hbdefinition}

\begin{hblemma}{thm:p15-w02b-p15-012-exponential-remainder}{proved=annex:p15-w02b-p15-012-exponential-remainder}{Feller's Lemma 1 (4.1)--(4.3)}
For \(n=1,2,\ldots\) put
\(\rho_n(t)=e^{it}-\sum_{k=0}^{n-1}(it)^k/k!\). Then
\(\rho_1(t)=i\int_0^te^{ix}\,dx\),
\(\rho_n(t)=i\int_0^t\rho_{n-1}(x)\,dx\) for \(n>1\), and
\[
 |\rho_n(t)|\leq\frac{|t|^n}{n!}\qquad\text{for every real }t .
 \tag{4.1}
\]
Feller states (4.1) for \(t>0\); the case \(t<0\) follows because
\(\rho_n(-t)\) is the complex conjugate of \(\rho_n(t)\).
\end{hblemma}

\begin{hblemma}{thm:p15-w02b-p15-012-xv4-derivatives}{proved=annex:p15-w02b-p15-012-xv4-derivatives}{Feller's Lemma 2 (4.5)}
If \(M_n<\infty\), the \(n\)th derivative of \(\varphi\) exists, is continuous,
and
\[
 \varphi^{(n)}(\zeta)=i^n\int_{-\infty}^{+\infty}e^{i\zeta x}x^n\,F\{dx\}.
 \tag{4.5}
\]
In particular \(\varphi^{(n)}(0)=i^nm_n\).
\end{hblemma}

\begin{hbtheorem}{thm:p15-w02b-p15-012-taylor-bound}{proved=annex:p15-w02b-p15-012-taylor-bound}{The Taylor bound (4.13)--(4.14)}
For real \(t,x,\zeta\) and \(n\geq1\),
\[
 \Bigl|e^{i\zeta x}\Bigl(e^{itx}-1-\frac{itx}{1!}-\cdots
        -\frac{(itx)^{n-1}}{(n-1)!}\Bigr)\Bigr|\leq\frac{|tx|^n}{n!}.
 \tag{4.13}
\]
If \(M_n<\infty\), then for all real \(\zeta\) and \(t\),
\[
 \Bigl|\varphi(\zeta+t)-\varphi(\zeta)-\frac t{1!}\varphi'(\zeta)-\cdots
     -\frac{t^{n-1}}{(n-1)!}\varphi^{(n-1)}(\zeta)\Bigr|
   \leq M_n\frac{|t|^n}{n!}.
 \tag{4.14}
\]
Feller prints a strict inequality in (4.14); the non-strict form is the one
that holds for every law, since both sides vanish for \(F=\delta_0\).
\end{hbtheorem}

\begin{hbtheorem}{thm:p15-w02b-p15-012-xv4-determinacy}{proved=annex:p15-w02b-p15-012-xv4-determinacy}{Feller's analytic moment-determinacy theorem (4.15)}
Suppose all moments of \(F\) exist and
\[
 \limsup_{n\to\infty}\frac1nM_n^{1/n}=\lambda<\infty .
 \tag{4.15}
\]
Then for every real \(\zeta\) and every real \(t\) with \(e\lambda|t|<1\)
(every real \(t\) when \(\lambda=0\)),
\[
 \varphi(\zeta+t)=\sum_{k=0}^\infty\frac{t^k}{k!}\varphi^{(k)}(\zeta).
\]
In particular this holds for \(|t|<1/(3\lambda)\), Feller's radius.
If \(G\) is any distribution with the same moments
\(m_n\), \(n\geq1\), then \(G=F\): \(F\) is uniquely determined by its moments.
\end{hbtheorem}

Feller's proof uses Stirling's formula to make the right side of (4.14)
vanish as \(n\to\infty\), concludes that \(\varphi\) is analytic near every
real point and hence determined by its power series at the origin, whose
coefficients are \(i^nm_n/n!\)
\cite[vol.~II, XV.4, p.~514]{Feller1971ProbabilityII}. The annex uses the
elementary bound \(n!\geq(n/e)^n\) in place of Stirling's formula, shows that
any competing law satisfies (4.15) with the same \(\lambda\), and carries out
the propagation from the origin along the real axis explicitly. It does not
use complex analytic continuation.

\begin{hbcorollary}{thm:p15-w02b-p15-012-xv4-bridge}{proved=annex:p15-w02b-p15-012-xv4-bridge}{The power-series criterion of VII.3 and VII.6 through XV.4}
Let a real law have all moments, and suppose
\[
 \sum_{n=0}^\infty\frac{m_{2n}s^{2n}}{(2n)!}<\infty
 \qquad\text{for some }s>0 .
\]
Then (4.15) holds with \(\lambda\leq1/s\), so the law is determined by its
moments. Equivalently, it suffices that \(\sum_n m_{2n}t^n/(2n)!\) converge
for some \(t>0\). For a law on \([0,\infty)\), convergence of
\(\sum_k\mu_k\lambda_0^k/k!\) for some \(\lambda_0>0\) implies this hypothesis
with \(s=\lambda_0\).
Condition (4.15) in turn implies Carleman's condition
\(\sum_nM_n^{-1/n}=\infty\), with the convention that a term with
\(M_n=0\) is \(+\infty\).
\end{hbcorollary}

Feller states in VII.3 that the book proves only the power-series criterion
and refers to VII.6 and XV.4 for it, and in VII.6 that the criterion holds for
laws not concentrated on \([0,\infty)\) by the argument of XV.4
\cite[vol.~II, VII.3, p.~228; VII.6, p.~234]{Feller1971ProbabilityII}.
The corollary supplies that link. On p.~515 Feller compares (4.15) with
Carleman's condition, describing his criterion as the weaker one
\cite[vol.~II, XV.4, p.~515]{Feller1971ProbabilityII}.

\begin{hbtheorem}{thm:p15-w02b-p15-012-xv4-examples}{proved=annex:p15-w02b-p15-012-xv4-examples}{Normal, exponential and lognormal boundaries of (4.15)}
\begin{enumerate}
\item For the standard normal law, (4.15) holds with \(\lambda=0\), so its
      characteristic function equals its Taylor series about every real
      point for every real increment.
\item For the exponential law of mean \(1\), \(M_n=n!\) and (4.15) holds with
      \(\lambda\leq1\).
\item For the lognormal moments \(M_n=e^{n^2/2}\) of
      \cref{thm:p15-feller-lognormal-moments}, \(M_n^{1/n}/n\to\infty\), so
      (4.15) fails, as it must, since those moments do not determine the law.
\end{enumerate}
\end{hbtheorem}

\section{Feller's continuity theorems for Laplace transforms}
\label{sec:p15-w02b-p15-012-laplace-continuity}

The inversion formula (6.4) shows that a law on \([0,\infty)\) is determined
by its Laplace transform. Feller's continuity theorem turns this uniqueness
into a convergence criterion: pointwise convergence of transforms is
equivalent to convergence of the laws, provided mass that escapes to infinity
is allowed in the limit \cite[vol.~II, XIII.1, pp.~429--433]{Feller1971ProbabilityII}.
The inputs are a sequence of laws, or more generally of measures that may have
infinite total mass, together with their transforms. The conclusions identify
the limit and say when no mass is lost.

\begin{hbdefinition}{def:p15-w02b-p15-012-vague}{Possibly defective distributions and Feller's convergence of measures}
A \emph{possibly defective distribution} on the line is a Borel measure \(F\)
with total mass \(F\{(-\infty,\infty)\}\leq1\); it is \emph{proper} when the
mass is \(1\) and \emph{defective} otherwise. More generally let \(U\) be a
Borel measure on the line with \(U\{I\}<\infty\) for every bounded interval
\(I\). A bounded interval is an \emph{interval of continuity} of \(U\) if
neither endpoint is an atom of \(U\). Measures \(U_n\) of the same kind
\emph{converge} to \(U\), written \(U_n\to U\), if
\[
 U_n\{I\}\longrightarrow U\{I\}\qquad\text{for every bounded interval of continuity }I\text{ of }U .
\]
This is the convergence Feller recalls from VIII.1 and VIII.6
\cite[vol.~II, XIII.1, p.~433]{Feller1971ProbabilityII}. A measure is
\emph{concentrated on} \([0,\infty)\) if it gives no mass to \((-\infty,0)\);
for it, \(U(x)=U\{[0,x]\}\) is its distribution function.
\end{hbdefinition}

\begin{hblemma}{thm:p15-w02b-p15-012-vague-tools}{proved=annex:p15-w02b-p15-012-vague-tools}{The convergence facts used by both continuity theorems}
Let \(U_n\) and \(U\) be measures finite on bounded intervals.
\begin{enumerate}
\item The atoms of \(U\) are countable, so its non-atoms are dense. If
      \(U_n\{(r,s]\}\to U\{(r,s]\}\) whenever \(r<s\) are non-atoms of \(U\),
      then \(U_n\to U\).
\item Let \(U_n\to U\), let \(a<b\) be non-atoms of \(U\), and let
      \(\mathcal G\) be a family of complex functions on \([a,b]\) with
      \(|g|\leq B\) for all \(g\in\mathcal G\) and uniformly equicontinuous:
      for every \(\epsilon>0\) there is \(\delta>0\) with
      \(|g(x)-g(y)|\leq\epsilon\) for all \(g\in\mathcal G\) whenever
      \(|x-y|\leq\delta\). Then
      \(\int_{(a,b]}g\,dU_n\to\int_{(a,b]}g\,dU\) uniformly in
      \(g\in\mathcal G\).
\item If \(U_n\to U\), every \(U_n\) and \(U\) have total mass at most \(1\),
      and \(g\) is continuous with \(g(x)\to0\) as \(|x|\to\infty\), then
      \(\int g\,dU_n\to\int g\,dU\).
\item If \(U_n\to U\) and \(U_n\to V\), then \(U=V\).
\item If every subsequence of \((U_n)\) has a further subsequence converging
      to \(U\), then \(U_n\to U\).
\item For probability distributions \(U_n\) and \(U\), \(U_n\to U\) holds if
      and only if \(U_n\Rightarrow U\) in the sense of
      \cref{def:p15-weak-law-convergence}.
\item Every sequence of probability distributions has a subsequence
      converging to a possibly defective distribution. If all the
      distributions are concentrated on \([0,\infty)\), so is the limit.
\end{enumerate}
\end{hblemma}

Item 7 is Feller's selection theorem, which he quotes from VIII.6
\cite[vol.~II, XIII.1, p.~431]{Feller1971ProbabilityII}; the annex derives it
from Durrett's selection construction (\cref{thm:p15-helly-selection}).

\begin{hbdefinition}{def:p15-w02b-p15-012-laplace-measure}{Laplace transforms of defective distributions and of measures (1.1), (1.5), (1.6)}
For a possibly defective distribution \(F\) concentrated on \([0,\infty)\),
the Laplace transform is
\[
 \varphi(\lambda)=\int_0^\infty e^{-\lambda x}\,F\{dx\},\qquad\lambda>0,
 \tag{1.1}
\]
the interval of integration being closed at \(0\). For a measure \(U\)
concentrated on \([0,\infty)\) and finite on bounded intervals, put
\[
 \omega(\lambda)=\int_0^\infty e^{-\lambda x}\,U\{dx\}.
 \tag{1.5}
\]
If \(\omega(a)<\infty\) then \(\omega(\lambda)<\infty\) for all \(\lambda>a\).
If the integral converges for \(\lambda>a\), the function \(\omega\) on
\((a,\infty)\) is the \emph{Laplace transform of} \(U\). When
\(U\{dx\}=u(x)\,dx\) with \(u\geq0\), (1.5) becomes the \emph{ordinary
Laplace transform} \(\omega(\lambda)=\int_0^\infty e^{-\lambda x}u(x)\,dx\)
of \(u\), formula (1.6). Here \(a\) may be any real number.
\end{hbdefinition}

\begin{hbtheorem}{thm:p15-w02b-p15-012-laplace-uniqueness}{proved=annex:p15-w02b-p15-012-laplace-uniqueness}{Feller's uniqueness theorems 1 and 1a and their corollary}
\begin{enumerate}
\item Two finite measures concentrated on \([0,\infty)\), in particular two
      possibly defective distributions, whose Laplace transforms agree on a
      half-line \(\lambda>\lambda_1\) are equal.
\item (Theorem~1a.) A measure \(U\) concentrated on \([0,\infty)\) whose
      Laplace transform (1.5) converges for \(\lambda>a\) is uniquely
      determined by the values of the transform on \(a<\lambda<\infty\).
\item (Corollary.) A nonnegative continuous function \(u\) on \([0,\infty)\)
      whose ordinary Laplace transform converges for \(\lambda>a\) is
      uniquely determined by that transform on \(a<\lambda<\infty\).
\end{enumerate}
\end{hbtheorem}

Feller proves Theorem~1 either by the inversion formula (6.4) or by the change
of variable \(y=e^{-x}\), which turns the transform at \(\lambda=1,2,\ldots\)
into the moments of a law on \([0,1]\), so that Hausdorff's uniqueness applies
\cite[vol.~II, XIII.1, p.~430]{Feller1971ProbabilityII}. The annex uses the
first route through \cref{thm:p15-w02b-p15-012-laplace-inversion}, extended
from probability distributions to finite measures by scaling. Theorem~1a
follows from Feller's \emph{translation principle}: if \(\omega(a')<\infty\),
then \(\omega(\lambda+a')\) is the transform of the finite measure
\(e^{-a'x}U\{dx\}\) \cite[vol.~II, XIII.1, (1.9), p.~432]{Feller1971ProbabilityII}.

\begin{hbtheorem}{thm:p15-w02b-p15-012-laplace-continuity}{proved=annex:p15-w02b-p15-012-laplace-continuity}{Feller's continuity theorem for Laplace transforms (XIII.1, Theorem~2)}
For \(n=1,2,\ldots\) let \(F_n\) be a probability distribution concentrated on
\([0,\infty)\), with transform \(\varphi_n\).
\begin{enumerate}
\item If \(F_n\to F\), where \(F\) is a possibly defective distribution with
      transform \(\varphi\), then \(\varphi_n(\lambda)\to\varphi(\lambda)\) for
      every \(\lambda>0\).
\item Conversely, if \(\varphi_n(\lambda)\) converges for each \(\lambda>0\) to
      a limit \(\varphi(\lambda)\), then \(\varphi\) is the transform of a
      possibly defective distribution \(F\) concentrated on \([0,\infty)\),
      and \(F_n\to F\).
\item The limit \(F\) is proper if and only if \(\varphi(\lambda)\to1\) as
      \(\lambda\to0+\). In that case \(F_n\Rightarrow F\).
\end{enumerate}
\end{hbtheorem}

Feller obtains the first part from his basic convergence theorem of VIII.1
and the second from the selection theorem of VIII.6, applied to subsequences
\cite[vol.~II, XIII.1, p.~431]{Feller1971ProbabilityII}. The annex follows
this route with \cref{thm:p15-w02b-p15-012-vague-tools} in place of those two
results.

\begin{hbtheorem}{thm:p15-w02b-p15-012-laplace-extended-continuity}{proved=annex:p15-w02b-p15-012-laplace-extended-continuity}{The extended continuity theorem (XIII.1, Theorem~2a)}
For \(n=1,2,\ldots\) let \(U_n\) be a measure concentrated on \([0,\infty)\)
whose transform \(\omega_n\) converges for \(\lambda>a\).
\begin{enumerate}
\item If \(\omega_n(\lambda)\to\omega(\lambda)<\infty\) for every
      \(\lambda>a\), then \(\omega\) is the Laplace transform of a measure
      \(U\) concentrated on \([0,\infty)\), and \(U_n\to U\).
\item Conversely, if \(U_n\to U\) and the sequence \(\{\omega_n(a)\}\) is
      bounded, then \(\omega_n(\lambda)\to\omega(\lambda)\) for
      \(\lambda>a\), where \(\omega\) is the transform of \(U\).
\end{enumerate}
\end{hbtheorem}

Feller reduces the first part to Theorem~2 by the translation principle:
for fixed \(\lambda_0>a\), the measures \(e^{-\lambda_0x}U_n\{dx\}/\omega_n(\lambda_0)\)
are probability distributions with transforms
\(\omega_n(\lambda+\lambda_0)/\omega_n(\lambda_0)\)
\cite[vol.~II, XIII.1, p.~433]{Feller1971ProbabilityII}. That reduction
divides by \(\omega(\lambda_0)\); the annex treats the remaining case
\(\omega(\lambda_0)=0\) separately, where the limit is the zero measure.

\begin{hbtheorem}{thm:p15-w02b-p15-012-laplace-continuity-examples}{proved=annex:p15-w02b-p15-012-laplace-continuity-examples}{A Poisson limit, escaping mass, and Feller's example (g)}
\begin{enumerate}
\item Let \(\mu>0\) and let \(F_n\) be the binomial law with \(n\) trials and
      success probability \(\mu/n\) (\(n>\mu\)). Then
      \[
       \varphi_n(\lambda)=\bigl(1-\tfrac\mu n(1-e^{-\lambda})\bigr)^n
        \longrightarrow\exp\{-\mu(1-e^{-\lambda})\}\qquad(\lambda>0),
      \]
      the limit is the transform of the Poisson law of mean \(\mu\) and tends
      to \(1\) as \(\lambda\to0+\), so \(F_n\Rightarrow\) Poisson. In
      particular \(\binom nk(\mu/n)^k(1-\mu/n)^{n-k}\to e^{-\mu}\mu^k/k!\)
      for every \(k\geq0\).
\item For \(F_n=\delta_n\), \(\varphi_n(\lambda)=e^{-n\lambda}\to0\), and
      \(F_n\to0\), the zero measure. For
      \(F_n=\tfrac12\delta_0+\tfrac12\delta_n\), \(\varphi_n\to\tfrac12\) and
      \(F_n\to\tfrac12\delta_0\). In both cases
      \(\lim_{\lambda\to0+}\varphi(\lambda)<1\), and the missing mass is the
      mass that escapes to infinity.
\item (Feller's example (g).) Let \(U_n\) attach the weight \(e^{n^2}\) to the
      point \(n\). Then \(U_n\to0\), but
      \(\omega_n(\lambda)=e^{n(n-\lambda)}\to\infty\) for every
      \(\lambda>0\). The boundedness of \(\{\omega_n(a)\}\) cannot be dropped
      from part~2 of \cref{thm:p15-w02b-p15-012-laplace-extended-continuity}.
\end{enumerate}
\end{hbtheorem}

Example (g) is Feller's \cite[vol.~II, XIII.1, p.~434]{Feller1971ProbabilityII}.
The binomial limit and the escaping-mass examples are added here as worked
instances of each part of the theorem.

\begin{hbtheorem}{thm:p15-w02b-p15-012-laplace-integer-values}{proved=annex:p15-w02b-p15-012-laplace-integer-values}{Feller's second uniqueness proof: the values at the integers}
Let \(F\) be a probability distribution concentrated on \([0,\infty)\), with
Laplace transform \(\varphi\), and let \(X\) have law \(F\). Put \(y=e^{-x}\)
and let \(G\) be the law of \(Y=e^{-X}\). Then \(G\) is a probability
distribution concentrated on \((0,1]\); for \(0<y<1\) and \(x=-\log y\),
\(G(y)=1-F(x)\) whenever \(x\) is a point of continuity of \(F\); and
\[
 \varphi(\lambda)=\int_0^\infty e^{-\lambda x}\,F\{dx\}=\int_0^1y^\lambda\,G\{dy\},
 \qquad\lambda>0.
 \tag{1.4}
\]
In particular \(\varphi(k)\) is the \(k\)th moment of \(G\), and the values
\(\varphi(1),\varphi(2),\ldots\) determine \(F\). The same conclusion holds for
a possibly defective \(F\) concentrated on \([0,\infty)\) if \(G\) is the image
of \(F\) under \(x\mapsto e^{-x}\) completed by an atom of mass
\(1-F\{[0,\infty)\}\) at \(y=0\).
\end{hbtheorem}

This is Feller's second proof of his Theorem~1: the change of variable
\(y=e^{-x}\) turns the transform at the positive integers into the moment
sequence of a law on \([0,1]\), and Hausdorff's uniqueness
(\cref{thm:p15-feller-hausdorff}) identifies that law
\cite[vol.~II, XIII.1, p.~430]{Feller1971ProbabilityII}. Feller justifies
(1.4) by matching Riemann sums; the annex uses the image-measure rule instead.
As he remarks, the conclusion is stronger than Theorem~1, since only countably
many values of the transform are used. The extension to defective \(F\) is
added here: the atom at \(0\) records the mass that \(F\) does not carry, and
it contributes nothing to the moments of positive order.

\begin{hbtheorem}{thm:p15-w02b-p15-012-laplace-xiii1-examples}{proved=annex:p15-w02b-p15-012-laplace-xiii1-examples}{Feller's transforms in XIII.1, examples (a)--(f)}
\begin{enumerate}
\item[(a)] If \(X\) takes the values \(0,1,\ldots\) with probabilities
      \(p_0,p_1,\ldots\) and generating function \(P(s)=\sum p_ns^n\), then
      \(\varphi(\lambda)=\sum p_ne^{-n\lambda}=P(e^{-\lambda})\) for
      \(\lambda>0\).
\item[(b)] For \(\alpha>0\) the gamma density
      \(f_\alpha(x)=(x^{\alpha-1}/\Gamma(\alpha))e^{-x}\), \(x>0\), has the
      transform
      \[
       \varphi_\alpha(\lambda)=\frac1{\Gamma(\alpha)}\int_0^\infty e^{-(\lambda+1)x}x^{\alpha-1}\,dx
         =\frac1{(\lambda+1)^\alpha},\qquad\lambda>0.
       \tag{1.3}
      \]
\item[(c)] If \(u(x)=x^a\) with \(a>-1\), then
      \(\omega(\lambda)=\Gamma(a+1)/\lambda^{a+1}\) for all \(\lambda>0\).
\item[(d)] If \(u(x)=e^{ax}\) with \(a>0\), then
      \(\omega(\lambda)=1/(\lambda-a)\) for \(\lambda>a\), and the integral
      (1.6) diverges for \(\lambda\leq a\).
\item[(e)] If \(u(x)=e^{x^2}\), the integral (1.6) diverges for every real
      \(\lambda\).
\item[(f)] For a possibly defective distribution \(F\) concentrated on
      \([0,\infty)\), with transform \(\varphi\),
      \[
       -\varphi'(\lambda)=\int_0^\infty e^{-\lambda x}x\,F\{dx\},\qquad\lambda>0.
       \tag{1.7}
      \]
      This is an integral of the form (1.5) with \(U\{dx\}=x\,F\{dx\}\), a
      measure finite on bounded intervals whose total mass is the first
      moment of \(F\) and may be infinite.
\end{enumerate}
\end{hbtheorem}

These are Feller's examples \cite[vol.~II, XIII.1, pp.~430--431]{Feller1971ProbabilityII}.
In (c) the exponent is printed as \(\alpha+1\) although the parameter is
named \(a\); the computation gives \(a+1\). Example (d) holds with the same
proof for every real \(a\); Feller states it for \(a>0\). Examples (d) and (e)
show the two ways in which the half-line of convergence in
\cref{def:p15-w02b-p15-012-laplace-measure} can fail to be all of
\(\lambda>0\): it may begin at a positive abscissa, or be empty. Example (f),
applied to a law of infinite mean such as \(F\{dx\}=(1+x)^{-2}dx\) on
\([0,\infty)\), gives a measure \(U\) of infinite total mass whose transform
is nonetheless finite for every \(\lambda>0\).

\section{Feller's continuity theorem for characteristic functions}
\label{sec:p15-w02b-p15-012-cf-continuity}

Feller's route to the continuity theorem for characteristic functions avoids
the inversion integral of \cref{thm:p15-characteristic-inversion-uniqueness}.
It smooths the law with a normal distribution, reads the smoothed density off
the characteristic function by a Parseval identity, and lets the smoothing
vanish \cite[vol.~II, XV.3, pp.~507--509]{Feller1971ProbabilityII}. The same
identity, applied along a convergent subsequence, gives the continuity theorem.
In this section \(\zeta\) is the transform variable, as in Feller's chapter XV.

\begin{hbdefinition}{def:p15-w02b-p15-012-normal-smoothing}{Normal densities and characteristic functions of defective distributions}
Feller writes \(\mathfrak n(x)=(2\pi)^{-1/2}e^{-\frac12x^2}\) for the standard
normal density. For \(a>0\), \(\mathfrak N_a\) denotes the normal distribution
with expectation \(0\) and variance \(a^2\), whose density is
\(a^{-1}\mathfrak n(x/a)\). The characteristic function of a possibly
defective distribution \(F\) on the line is
\(\varphi(\zeta)=\int e^{i\zeta x}\,F\{dx\}\); thus
\(\varphi(0)=F\{(-\infty,\infty)\}\) and \(|\varphi|\leq\varphi(0)\leq1\).
A sequence \(F_n\) \emph{converges properly} to \(F\) if \(F_n\to F\) and
\(F\) is proper.
\end{hbdefinition}

\begin{hblemma}{thm:p15-w02b-p15-012-parseval}{proved=annex:p15-w02b-p15-012-parseval}{The Parseval relation and normal smoothing (3.1)--(3.4)}
Let \(F\) be a possibly defective distribution with characteristic function
\(\varphi\), and let \(G\) be a probability distribution with characteristic
function \(\gamma\). For every real \(t\),
\[
 e^{-i\zeta t}\varphi(\zeta)=\int_{-\infty}^{+\infty}e^{i\zeta(x-t)}\,F\{dx\},
 \tag{3.1}
\]
\[
 \int_{-\infty}^{+\infty}e^{-i\zeta t}\varphi(\zeta)\,G\{d\zeta\}
   =\int_{-\infty}^{+\infty}\gamma(x-t)\,F\{dx\}.
 \tag{3.2}
\]
For \(a>0\) and the normal distribution with density \(a\mathfrak n(a\zeta)\),
whose characteristic function is \(\sqrt{2\pi}\,\mathfrak n(\zeta/a)\),
\[
 \int_{-\infty}^{+\infty}e^{-i\zeta t}\varphi(\zeta)\,a\mathfrak n(a\zeta)\,d\zeta
   =\sqrt{2\pi}\int_{-\infty}^{+\infty}\mathfrak n\Bigl(\frac{x-t}a\Bigr)F\{dx\},
 \tag{3.3}
\]
which is the same as
\[
 \frac1{2\pi}\int_{-\infty}^{+\infty}e^{-i\zeta t}\varphi(\zeta)e^{-\frac12a^2\zeta^2}\,d\zeta
   =\int_{-\infty}^{+\infty}\frac1a\,\mathfrak n\Bigl(\frac{t-x}a\Bigr)F\{dx\}
   =:f_a(t).
 \tag{3.4}
\]
For every bounded continuous \(g\),
\(\int g(t)f_a(t)\,dt\to\int g\,dF\) as \(a\to0+\). If \(F\) is proper,
\(f_a\) is the density of a probability distribution
\(\mathfrak N_a\star F\), and \(\mathfrak N_a\star F\Rightarrow F\) as
\(a\to0+\).
\end{hblemma}

Feller prints the constant on the left of (3.4) as \(1/\sqrt{2\pi}\); dividing
(3.3) by \(a\sqrt{2\pi}\) gives \(1/(2\pi)\), the constant used here
\cite[vol.~II, XV.3, p.~507]{Feller1971ProbabilityII}. With \(1/(2\pi)\), the
case \(t=0\), \(F=\delta_0\) reads \(\frac1{2\pi}\cdot\sqrt{2\pi}/a=\mathfrak n(0)/a\), as it must.

\begin{hbtheorem}{thm:p15-w02b-p15-012-cf-uniqueness}{proved=annex:p15-w02b-p15-012-cf-uniqueness}{Uniqueness by normal smoothing (XV.3, Theorem~1)}
Distinct probability distributions have distinct characteristic functions.
More generally, let \(\varphi\) be any bounded measurable function on the
line. At most one possibly defective distribution \(F\) satisfies (3.4) with
this \(\varphi\) on the left for all \(a>0\) and all real \(t\).
\end{hbtheorem}

The second sentence is the form Feller uses in the next proof, where the limit
function \(\varphi\) is not yet known to be a characteristic function
\cite[vol.~II, XV.3, p.~508]{Feller1971ProbabilityII}.

\begin{hblemma}{thm:p15-w02b-p15-012-cf-lemma}{proved=annex:p15-w02b-p15-012-cf-lemma}{Feller's convergence lemma (XV.3)}
Let \(F_n\) be a probability distribution with characteristic function
\(\varphi_n\). If \(\varphi_n(\zeta)\to\varphi(\zeta)\) for all \(\zeta\), then
there exists a possibly defective distribution \(F\) such that \(F_n\to F\),
and \(F\) satisfies (3.4) with the limit function \(\varphi\) on the left.
\end{hblemma}

\begin{hbtheorem}{thm:p15-w02b-p15-012-cf-continuity}{proved=annex:p15-w02b-p15-012-cf-continuity}{Feller's continuity theorem (XV.3, Theorem~2) and its corollary}
In order that a sequence \(\{F_n\}\) of probability distributions converge
properly to a probability distribution \(F\), it is necessary and sufficient
that the characteristic functions \(\varphi_n\) converge pointwise to a limit
\(\varphi\) that is continuous in some neighbourhood of the origin. In this
case \(\varphi\) is the characteristic function of \(F\), \(\varphi\) is
continuous everywhere, and \(\varphi_n\to\varphi\) uniformly in every finite
interval.

(Corollary.) A continuous function which is the pointwise limit of a sequence
of characteristic functions is itself a characteristic function.
\end{hbtheorem}

Feller's sufficiency argument lets \(a\to\infty\) in (3.3): the left side is
an average under a normal law of variance \(a^{-2}\), which concentrates at the
origin, and the right side is at most the total mass of \(F\)
\cite[vol.~II, XV.3, p.~509]{Feller1971ProbabilityII}. That step uses
continuity of \(\varphi\) at the origin only. For necessity Feller cites the
corollary of VIII.1; the annex proves the uniform convergence directly from
tightness. Durrett's version of the same theorem, with a tightness estimate in
place of normal smoothing, is \cref{thm:p15-characteristic-continuity-theorem}.

\begin{hbtheorem}{thm:p15-w02b-p15-012-cf-examples}{proved=annex:p15-w02b-p15-012-cf-examples}{A uniform limit and a discontinuous limit}
\begin{enumerate}
\item Let \(F_n\) put mass \(1/n\) at each point \(k/n\), \(k=1,\ldots,n\).
      Then \(|\varphi_n(\zeta)-\varphi(\zeta)|\leq|\zeta|/n\), where
      \(\varphi(\zeta)=\int_0^1e^{i\zeta x}\,dx\) is the characteristic
      function of the uniform law on \([0,1]\). Hence \(F_n\) converges
      properly to that uniform law.
\item Let \(F_n=\tfrac12\delta_0+\tfrac12\mathfrak N_{\sqrt n}\). Then
      \(\varphi_n(\zeta)=\tfrac12+\tfrac12e^{-n\zeta^2/2}\) converges to
      \(\varphi\) with \(\varphi(0)=1\) and \(\varphi(\zeta)=\tfrac12\) for
      \(\zeta\neq0\), a limit discontinuous at the origin, and
      \(F_n\to\tfrac12\delta_0\), a defective distribution whose
      characteristic function is the constant \(\tfrac12\), not \(\varphi\).
\end{enumerate}
\end{hbtheorem}

The second example shows both conclusions of
\cref{thm:p15-w02b-p15-012-cf-continuity} failing when the limit function is
discontinuous at the origin: mass escapes, and the pointwise limit is not the
characteristic function of the limit.

\begin{hbtheorem}{thm:p15-w02b-p15-012-convolution-powers}{proved=annex:p15-w02b-p15-012-convolution-powers}{Convolution powers escape to infinity (XV.3, example (a))}
Let \(U\) be a probability distribution whose characteristic function
\(\omega\) is real and nonnegative, and let \(U^{n\star}\) be its
\(n\)-fold convolution, with characteristic function \(\omega^n\). If \(U\) is
not concentrated at the origin, then \(U^{n\star}\to0\), the zero measure: the
mass of every bounded interval tends to \(0\). For every probability
distribution \(G\) not concentrated at the origin, likewise \(G^{n\star}\to0\).
\end{hbtheorem}

Feller identifies the set where \(\omega=1\) through his classification of
the points where \(|\varphi|=1\) (lemma~4 of XV.1), concludes from (3.4) that
the limit has vanishing smoothed densities, and passes to general \(G\) by
symmetrization \cite[vol.~II, XV.3, p.~508]{Feller1971ProbabilityII}. The
annex needs less than that lemma: it shows directly that the set
\(\{\omega=1\}\) has Lebesgue measure zero. For \(G\) not concentrated at a
single point it uses the symmetrized law; for \(G=\delta_b\) with \(b\neq0\)
the conclusion is immediate, since \(\delta_{nb}\) leaves every bounded
interval.

\begin{hbtheorem}{thm:p15-w02b-p15-012-polya-feller}{proved=annex:p15-w02b-p15-012-polya-feller}{P\'olya's criterion by inscribed polygons (XV.3, example (b))}
Let \(\omega\) be a real even function with \(\omega(0)=1\), \(\omega\geq0\),
continuous at \(0\), and convex and nonincreasing on \([0,\infty)\). Then
\(\omega\) is a characteristic function. The three added hypotheses cannot be
dropped from Feller's printed statement, which asks only that \(\omega\) be
real, even, equal to \(1\) at \(0\), and have a convex graph on \([0,\infty)\):
\begin{enumerate}
\item \(\omega(\zeta)=1+\zeta^2\) is nonnegative, continuous and convex, but
      not bounded by \(1\);
\item \(\omega(\zeta)=-\tfrac12+\tfrac32e^{-|\zeta|}\) is continuous, convex
      and nonincreasing on \([0,\infty)\), but negative for large
      \(|\zeta|\);
\item \(\omega=\mathbf 1_{\{0\}}\) is nonnegative, nonincreasing and convex on
      \([0,\infty)\), but discontinuous at \(0\);
\end{enumerate}
and none of the three is a characteristic function.
\end{hbtheorem}

Feller's argument has two steps \cite[vol.~II, XV.2, example (b),
pp.~504--505; XV.3, example (b), p.~509]{Feller1971ProbabilityII}. An even
function \(\omega\geq0\) with \(\omega(0)=1\) whose graph on \([0,\infty)\) is a
convex polygon with finitely many sides, the last of them on the axis, is a
mixture \(\sum p_k(1-|\zeta|/a_k)^+\) of rescaled copies of the
characteristic function (2.5), and hence a characteristic function. A general \(\omega\) is the pointwise
limit of its inscribed polygons, which are of that kind, and the corollary
of \cref{thm:p15-w02b-p15-012-cf-continuity} applies. Feller's hypotheses
for the polygon case include \(\omega\geq0\) and a last side on the axis,
which makes the polygon nonincreasing; the three examples show that
nonnegativity, monotonicity and continuity at the origin must also be required
in the general statement. The polygon representation is the finite-piece case of
\cref{thm:p15-criteria-polya}, Durrett's form of the criterion, which also
treats a positive limit at infinity.

\begin{hbtheorem}{thm:p15-w02b-p15-012-fourier-inversion}{proved=annex:p15-w02b-p15-012-fourier-inversion}{Fourier inversion by normal smoothing (XV.3, Theorem~3 and corollary)}
Write \(\varphi\in L\) if \(|\varphi|\) is integrable over
\((-\infty,\infty)\). Let \(\varphi\) be the characteristic function of the
probability distribution \(F\), and suppose \(\varphi\in L\). Then \(F\) has a
bounded continuous density
\[
 f(x)=\frac1{2\pi}\int_{-\infty}^{+\infty}e^{-i\zeta x}\varphi(\zeta)\,d\zeta.
 \tag{3.5}
\]
(Corollary.) If \(\varphi\geq0\), then \(\varphi\in L\) if and only if \(F\)
has a bounded density. If \(F\) has a density \(f\leq M\), then
\[
 \frac1{2\pi}\int_{-\infty}^{+\infty}\varphi(\zeta)e^{-\frac12a^2\zeta^2}\,d\zeta
  =\frac1{\sqrt{2\pi}\,a}\int_{-\infty}^{+\infty}e^{-x^2/(2a^2)}f(x)\,dx\leq M
 \tag{3.7}
\]
for every \(a>0\), and \(\int\varphi\,d\zeta\leq2\pi M\).
\end{hbtheorem}

The statement of Theorem~3 is that of
\cref{thm:p15-characteristic-integrable-inversion}, which Durrett proves from
the inversion integral over bounded intervals. Feller's proof is different: the
smoothed densities \(f_a\) of (3.4) converge boundedly to the right side of
(3.5), and the smoothed laws converge to \(F\)
\cite[vol.~II, XV.3, pp.~509--510]{Feller1971ProbabilityII}.

\begin{hbtheorem}{thm:p15-w02b-p15-012-xv3-inversion-examples}{proved=annex:p15-w02b-p15-012-xv3-inversion-examples}{Plancherel, densities and distribution functions (XV.3, examples (c)--(e))}
\begin{enumerate}
\item[(c)] (Plancherel identity.) Let \(F\) have a density \(f\) and
      characteristic function \(\varphi\). Let \({}^0F\) be the symmetrized
      distribution, the law of \(X-X'\) for independent \(X,X'\) with law
      \(F\); its characteristic function is \(|\varphi|^2\), and it has the
      density
      \[
       {}^0f(x)=\int_{-\infty}^{+\infty}f(y+x)f(y)\,dy .
       \tag{3.9}
      \]
      Then \(|\varphi|^2\in L\) if and only if \(f^2\in L\), and in that case
      \[
       \int_{-\infty}^{+\infty}f^2(y)\,dy=\frac1{2\pi}\int_{-\infty}^{+\infty}|\varphi(\zeta)|^2\,d\zeta .
       \tag{3.8}
      \]
\item[(d)] (Continuity theorem for densities.) Let \(\varphi_n\) and
      \(\varphi\) be integrable characteristic functions with
      \[
       \int_{-\infty}^{+\infty}|\varphi_n(\zeta)-\varphi(\zeta)|\,d\zeta\longrightarrow0 .
       \tag{3.10}
      \]
      Then the bounded continuous densities \(f_n\) and \(f\) given by (3.5)
      satisfy
      \(|f_n(x)-f(x)|\leq(2\pi)^{-1}\int|\varphi_n(\zeta)-\varphi(\zeta)|\,d\zeta\)
      for all \(x\), so \(f_n\to f\) uniformly.
\item[(e)] (Inversion for distribution functions.) Let \(F\) have
      characteristic function \(\varphi\) and let \(h>0\). If
      \(\zeta\mapsto\varphi(\zeta)(1-e^{-i\zeta h})/(i\zeta h)\) is integrable,
      in particular if \(|\varphi(\zeta)|=O(1/|\zeta|)\) as
      \(|\zeta|\to\infty\), then for every real \(x\)
      \[
       \frac{F(x+h)-F(x)}h=\frac1{2\pi}\int_{-\infty}^{+\infty}\varphi(\zeta)\,
         \frac{1-e^{-i\zeta h}}{i\zeta h}\,e^{-i\zeta x}\,d\zeta .
       \tag{3.11}
      \]
      Without any hypothesis on \(\varphi\), let \(F_a=\mathfrak N_a\star F\).
      Then for all real \(x\)
      \[
       \frac{F_a(x+h)-F_a(x)}h=\frac1{2\pi}\int_{-\infty}^{+\infty}\varphi(\zeta)
         e^{-\frac12a^2\zeta^2}\,\frac{1-e^{-i\zeta h}}{i\zeta h}\,e^{-i\zeta x}\,d\zeta,
       \tag{3.12}
      \]
      and if \(x\) and \(x+h\) are points of continuity of \(F\), the right
      side tends to \((F(x+h)-F(x))/h\) as \(a\to0+\).
\end{enumerate}
At \(\zeta=0\) the factor \((1-e^{-i\zeta h})/(i\zeta h)\) is given its
limiting value \(1\).
\end{hbtheorem}

These are Feller's examples and his note on inversion formulas
\cite[vol.~II, XV.3, pp.~510--511]{Feller1971ProbabilityII}. In (3.11) the
left side is the density of the convolution of \(F\) with the uniform
distribution on \([-h,0]\), whose characteristic function is the extra factor;
so (3.11) is (3.5) for that convolution. Formula (3.12) is the same identity
for the smoothed law \(F_a\), and it needs no integrability hypothesis. The
inversion of \cref{thm:p15-characteristic-inversion-uniqueness} is the
traditional variant that replaces the normal factor by truncation of the
integral to \([-T,T]\) and lets \(T\to\infty\).

\begin{hbtheorem}{thm:p15-w02b-p15-012-lattice-inversion}{proved=annex:p15-w02b-p15-012-lattice-inversion}{Lattice inversion (XV.3, Theorem~4)}
Let \(h>0\), let \(b\) be real, and let \(F\) attribute weight \(p_k\geq0\) to
the point \(b+kh\), \(k\in\mathbb Z\), with \(\sum p_k=1\), so that
\[
 \varphi(\zeta)=\sum_{k=-\infty}^{+\infty}p_ke^{i(b+kh)\zeta}.
 \tag{3.13}
\]
Then for every integer \(r\)
\[
 p_r=\frac h{2\pi}\int_{-\pi/h}^{\pi/h}\varphi(\zeta)e^{-i(b+rh)\zeta}\,d\zeta .
 \tag{3.14}
\]
\end{hbtheorem}

This is the lattice formula of \cref{thm:p15-characteristic-atom-inversion}
with \(x=b+rh\); Feller's proof integrates the series (3.13) term by term
\cite[vol.~II, XV.3, p.~511]{Feller1971ProbabilityII}.

\section{The exponential formula of semigroup theory}
\label{sec:p15-w02b-p15-012-semigroups}

In VII.5 Feller notes that the Poisson randomization behind (5.6) applies to
any continuous semigroup of bounded operators, and defers the proof to X.9
\cite[vol.~II, VII.5, p.~232]{Feller1971ProbabilityII}. This section gives
that general theorem in Feller's setting of X.8. Its input is a continuous
semigroup of contractions on a space of bounded functions; its conclusion is
that the semigroup is the strong limit of the explicit series built from the
single operator \(\mathfrak Q(h)\) as \(h\to0+\).

\begin{hbdefinition}{def:p15-w02b-p15-012-semigroup}{Function spaces, contractions and continuous semigroups}
Let \(\Sigma\) be a set and let \(\mathscr L\) be a family of bounded real
functions on \(\Sigma\) with the norm \(\|u\|=\sup_x|u(x)|\), such that
(i) every linear combination of members of \(\mathscr L\) is in
\(\mathscr L\); (ii) if \(u_n\in\mathscr L\) and \(\|u_n-u\|\to0\) then
\(u\in\mathscr L\); (iii) if \(u\in\mathscr L\) then \(u^+\) and \(u^-\)
are in \(\mathscr L\). A linear map \(T\) from \(\mathscr L\) to
\(\mathscr L\) is an \emph{endomorphism} if \(\|Tu\|\leq m\|u\|\) for a
constant \(m\); the least such \(m\) is \(\|T\|\). It is \emph{positive} if
\(u\geq0\) implies \(Tu\geq0\). A \emph{contraction} is a positive
endomorphism with \(\|T\|\leq1\). If the constant \(1\) is in \(\mathscr L\),
a positive \(T\) with \(T1=1\) is a \emph{transition operator}.
Endomorphisms \(T_n\) \emph{converge} to \(T\), written \(T_n\to T\), if
\(\|T_nu-Tu\|\to0\) for each \(u\in\mathscr L\). A family
\(\{\mathfrak Q(t):t\geq0\}\) of endomorphisms is a \emph{semigroup} if
\(\mathfrak Q(s+t)=\mathfrak Q(s)\mathfrak Q(t)\) for \(s,t>0\); a semigroup
of contractions is \emph{continuous} if \(\mathfrak Q(0)=\mathbf 1\), the
identity, and \(\mathfrak Q(h)\to\mathbf 1\) as \(h\to0+\).
\end{hbdefinition}

These are Feller's conventions \cite[vol.~II, X.8, pp.~350--352]{Feller1971ProbabilityII}.
His typical cases are the bounded continuous functions vanishing at infinity
on the line, the continuous functions on a closed bounded interval, and the
bounded sequences.

\begin{hblemma}{thm:p15-w02b-p15-012-semigroup-continuity}{proved=annex:p15-w02b-p15-012-semigroup-continuity}{Completeness and the continuity estimate (8.7)}
The space \(\mathscr L\) is complete in its norm. If \(\{\mathfrak Q(t)\}\) is
a continuous semigroup of contractions, then for \(0\leq t'<t''\) and
\(u\in\mathscr L\)
\[
 \|\mathfrak Q(t'')u-\mathfrak Q(t')u\|\leq\|\mathfrak Q(t''-t')u-u\|,
 \tag{8.7}
\]
so \(t\mapsto\mathfrak Q(t)u\) is uniformly continuous from \([0,\infty)\) to
\(\mathscr L\). The operators of the semigroup commute.
\end{hblemma}

\begin{hbtheorem}{thm:p15-w02b-p15-012-pseudo-poisson}{proved=annex:p15-w02b-p15-012-pseudo-poisson}{Pseudo-Poisson semigroups (X.9, Theorem~1)}
Let \(T\) be a contraction on \(\mathscr L\) and \(\alpha>0\). For
\(t\geq0\) the series
\[
 \mathfrak Q(t)=e^{-\alpha t}\sum_{n=0}^\infty\frac{(\alpha t)^n}{n!}T^n,
 \tag{9.1}
\]
abbreviated \(\mathfrak Q(t)=e^{\alpha t(T-\mathbf 1)}\) (9.2), converges in
operator norm, and the operators form a continuous semigroup of contractions
with \(\|\mathfrak Q(t)\|\leq e^{-\alpha t+\alpha t\|T\|}\leq1\). If \(T\) is a
transition operator, so is every \(\mathfrak Q(t)\). Such semigroups are
called \emph{pseudo-Poissonian}, generated by \(\alpha(T-\mathbf 1)\).
\end{hbtheorem}

\begin{hbtheorem}{thm:p15-w02b-p15-012-semigroup-exponential}{proved=annex:p15-w02b-p15-012-semigroup-exponential}{The exponential formula (X.9, Theorem~2)}
Let \(\{\mathfrak Q(t)\}\) be a continuous semigroup of contractions on
\(\mathscr L\). For \(h>0\) and \(t\geq0\) put
\[
 \mathfrak Q_h(t)=e^{-t/h}\sum_{n=0}^\infty\frac{(t/h)^n}{n!}\,\mathfrak Q(nh).
 \tag{9.3}
\]
Then \(\{\mathfrak Q_h(t)\}\) is the pseudo-Poisson semigroup generated by the
endomorphism \(h^{-1}[\mathfrak Q(h)-\mathbf 1]\), and
\[
 \mathfrak Q_h(t)\longrightarrow\mathfrak Q(t),\qquad h\to0 .
 \tag{9.4}
\]
Quantitatively, if \(\epsilon>0\), \(u\in\mathscr L\) and \(\delta>0\) satisfy
\(\|\mathfrak Q(s)u-u\|\leq\epsilon\) for \(0\leq s<\delta\), then
\[
 \|\mathfrak Q_h(t)u-\mathfrak Q(t)u\|\leq\epsilon+2\|u\|\,th\delta^{-2},
\]
so (9.4) holds uniformly in finite \(t\)-intervals.
\end{hbtheorem}

Feller's proof starts from the identity (9.5), which writes the difference as
a Poisson average of \(\mathfrak Q(nh)u-\mathfrak Q(t)u\), and splits the
average by Chebyshev's inequality
\cite[vol.~II, X.9, p.~354]{Feller1971ProbabilityII}. His (9.6) prints the
range of the good terms as \(|nh-t|<\eta t\); the estimate that the proof uses,
and that the annex proves, is for \(|nh-t|<\delta\).
\Cref{thm:p15-w02b-p15-012-exponential-formula} is the special case of the
translation semigroup, where the Poisson average is (5.7).

\begin{hbtheorem}{thm:p15-w02b-p15-012-subordinated-variant}{proved=annex:p15-w02b-p15-012-subordinated-variant}{Feller's variant with compound Poisson weights (9.7)--(9.11)}
Let \(p_0,p_1,\ldots\geq0\) with \(\sum p_j=1\), put
\(p(\zeta)=p_0+p_1\zeta+\cdots\) (9.8), assume
\(m_1=\sum_jjp_j>0\) and \(m_2=\sum_jj^2p_j<\infty\), and let \(\alpha>0\). For
\(s\geq0\) define \(u_n(s)\) by
\[
 \sum_{n=0}^\infty u_n(s)\zeta^n=\exp\bigl(s\alpha[p(\zeta)-1]\bigr),\qquad|\zeta|\leq1 .
 \tag{9.7}
\]
Then \(\{u_n(s)\}\) is a probability distribution with expectation \(bs\) and
variance \(cs\), where \(b=\alpha m_1\) and \(c=\alpha m_2\). For a continuous
semigroup of contractions and \(h>0\), the operators
\[
 \mathfrak Q_h(t)=\sum_nu_n\Bigl(\frac t{bh}\Bigr)\mathfrak Q(nh)
   =\sum_nu_n\Bigl(\frac t{bh}\Bigr)\mathfrak Q^n(h)
 \tag{9.9}
\]
satisfy \(\mathfrak Q_h(t)\to\mathfrak Q(t)\) as \(h\to0\), uniformly in finite
\(t\)-intervals, with
\(\|\mathfrak Q_h(t)u-\mathfrak Q(t)u\|\leq\epsilon+2\|u\|\,cth/(b\delta^2)\)
for \(\epsilon,\delta\) as in \cref{thm:p15-w02b-p15-012-semigroup-exponential}.
Moreover, with the contraction
\[
 \mathfrak Q^\#=p(\mathfrak Q(h))=\sum_{n=0}^\infty p_n\mathfrak Q^n(h)
   =\sum_{n=0}^\infty p_n\mathfrak Q(nh),
 \tag{9.10}
\]
the family \(\{\mathfrak Q_h(t)\}_{t\geq0}\) is the pseudo-Poisson semigroup
\[
 \mathfrak Q_h(t)=\exp\Bigl(\frac{\alpha t}{bh}(\mathfrak Q^\#-\mathbf 1)\Bigr)
 \tag{9.11}
\]
in the sense of (9.1) and (9.2).
\end{hbtheorem}

Feller states the variant in small print, takes the weights to be an
infinitely divisible distribution of the form (9.7), and asserts (9.11)
formally, printing it without the factor \(t\)
\cite[vol.~II, X.9, p.~355]{Feller1971ProbabilityII}. The theorem above takes
(9.7) as the definition of the weights, adds the moment assumptions under
which Feller's expectation \(bt\) and variance \(ct\) exist, and proves (9.11)
by rearranging absolutely convergent series. Feller credits the observation to
K.~L.~Chung \cite[vol.~II, X.9, p.~355, n.~32]{Feller1971ProbabilityII}.

\begin{hbtheorem}{thm:p15-w02b-p15-012-semigroup-examples}{proved=annex:p15-w02b-p15-012-semigroup-examples}{Translations on uniformly continuous functions, and the limits of strong convergence}
Let \(\mathscr L_u\) be the bounded uniformly continuous real functions on
\([0,\infty)\) and \(T(\theta)u(x)=u(x+\theta)\).
\begin{enumerate}
\item \(\mathscr L_u\) has properties (i)--(iii), each \(T(\theta)\) is a
      transition operator on it, and \(\{T(\theta)\}\) is a continuous
      semigroup of contractions.
\item For this semigroup \(\mathfrak Q_h(\theta)\) of (9.3) is Feller's
      interpolation operator \(A_h(\theta)\) of (5.7), and
      \cref{thm:p15-w02b-p15-012-semigroup-exponential} gives
      \(\|A_h(\theta)u-T(\theta)u\|\to0\) uniformly for \(\theta\) in finite
      intervals, for every \(u\in\mathscr L_u\).
\item The convergence is not in operator norm: on \(\mathscr L_u\),
      \(\|A_h(\theta)-T(\theta)\|=2\) whenever \(\theta>0\) and \(\theta/h\)
      is not an integer.
\end{enumerate}
\end{hbtheorem}

\section{Laplace transforms of semigroups and the abstract inversion formula}
\label{sec:p15-w02b-p15-012-resolvents}

Feller states that an abstract version of the Laplace-integral inversion (6.6)
appears in XIII.9 \cite[vol.~II, VII.6, p.~233]{Feller1971ProbabilityII}.
Section XIII.9 constructs the Laplace transform of a semigroup, the resolvent;
the inversion formula itself is Theorem~3 of XIII.10, (10.24), whose proof
Feller identifies with the gamma-kernel argument of VII.1
\cite[vol.~II, XIII.9, pp.~454--457; XIII.10, pp.~462--463]{Feller1971ProbabilityII}.
This section gives the construction, the resolvent limit (9.13), the
uniqueness lemma, the derivative formula (10.25) and the inversion (10.24).
Throughout, \(\mathscr L\) and \(\{\mathfrak Q(t)\}\) are as in
\cref{def:p15-w02b-p15-012-semigroup}, the semigroup being a continuous
semigroup of contractions.

\begin{hbdefinition}{def:p15-w02b-p15-012-resolvent}{Integrals of a semigroup and the resolvent (9.1), (9.5), (9.6)}
For a finite measure \(G\) on \([0,\infty)\) and \(w\in\mathscr L\), the
function \(\int_0^\infty\mathfrak Q(s)w\,G\{ds\}\) is defined at each
\(x\in\Sigma\) as the ordinary integral of the bounded continuous numerical
function \(s\mapsto\mathfrak Q(s)w(x)\). When \(G=F\) is a probability
distribution the resulting operator is written
\(E=\int_0^\infty\mathfrak Q(s)\,F\{ds\}\) (9.1), and \(Ew\) as in (9.5). For
\(G\{ds\}=e^{-\lambda s}\,ds\) with \(\lambda>0\) the operator is the
\emph{resolvent}
\[
 \mathfrak R(\lambda)=\int_0^\infty e^{-\lambda s}\mathfrak Q(s)\,ds .
 \tag{9.6}
\]
For fixed \(w\), derivatives \(\mathfrak R^{(n)}(\lambda)w\) are taken in the
norm of \(\mathscr L\).
\end{hbdefinition}

\begin{hblemma}{thm:p15-w02b-p15-012-semigroup-integral}{proved=annex:p15-w02b-p15-012-semigroup-integral}{Existence and the commutation rule (9.2)}
Let \(G\) be a finite measure on \([0,\infty)\) and \(w\in\mathscr L\).
\begin{enumerate}
\item \(\int\mathfrak Q(s)w\,G\{ds\}\) belongs to \(\mathscr L\). It is the
      norm limit, as \(\eta\to0+\), of the Riemann sums
      \(\sum_k\mathfrak Q(t_k)w\,G\{[k\eta,(k+1)\eta)\}\) with
      \(t_k\in[k\eta,(k+1)\eta)\).
\item The map \(w\mapsto\int\mathfrak Q(s)w\,G\{ds\}\) is a positive
      endomorphism of norm at most \(G\{[0,\infty)\}\); for a probability
      distribution \(F\), \(E\) is a contraction, and \(E1=1\) when every
      \(\mathfrak Q(t)\) is a transition operator.
\item For \(t\geq0\),
      \(\mathfrak Q(t)E=E\mathfrak Q(t)=\int_0^\infty\mathfrak Q(t+s)\,F\{ds\}\)
      (9.2), and the same holds for every finite \(G\).
\end{enumerate}
\end{hblemma}

Feller defines (9.1) by Riemann sums, as a special case of a Bochner integral
\cite[vol.~II, XIII.9, p.~455]{Feller1971ProbabilityII}. The lemma defines it
pointwise, which is the reading Feller gives of \(Ew(x)\), and then shows that
the Riemann sums converge to it in norm.

\begin{hbtheorem}{thm:p15-w02b-p15-012-resolvent-limit}{proved=annex:p15-w02b-p15-012-resolvent-limit}{The resolvent: contraction, limit (9.13) and uniqueness}
\begin{enumerate}
\item For \(\lambda>0\), \(\lambda\mathfrak R(\lambda)\) is a contraction that
      commutes with every \(\mathfrak Q(t)\); it is a transition operator when
      every \(\mathfrak Q(t)\) is.
\item (Theorem~2 of XIII.9.) As \(\lambda\to\infty\),
      \(\lambda\mathfrak R(\lambda)\to\mathbf 1\) (9.13), with
      \[
       \|\lambda\mathfrak R(\lambda)w-w\|\leq
         \int_0^\infty\|\mathfrak Q(t)w-w\|\,\lambda e^{-\lambda t}\,dt .
       \tag{9.14}
      \]
\item (Lemma of XIII.9.) The values \(\mathfrak R(\lambda)w\) for all
      \(\lambda>0\) and \(w\in\mathscr L\) determine the semigroup.
\end{enumerate}
\end{hbtheorem}

For the uniqueness lemma Feller applies the corollary of XIII.1 to the
numerical function \(t\mapsto\mathfrak Q(t)w(x)\). That corollary, as stated in
\cref{thm:p15-w02b-p15-012-laplace-uniqueness}, concerns nonnegative functions,
while \(\mathfrak Q(t)w(x)\) may change sign; the annex therefore uses (6.6),
which recovers any bounded continuous function from its Laplace integral, as
Feller himself points out after the corollary
\cite[vol.~II, XIII.1, p.~433]{Feller1971ProbabilityII}.

\begin{hbtheorem}{thm:p15-w02b-p15-012-resolvent-derivatives}{proved=annex:p15-w02b-p15-012-resolvent-derivatives}{Derivatives of the resolvent (10.25)}
For every \(w\in\mathscr L\), the map \(\lambda\mapsto\mathfrak R(\lambda)w\)
from \((0,\infty)\) to \(\mathscr L\) has norm derivatives of all orders, and
\[
 (-1)^n\mathfrak R^{(n)}(\lambda)=\int_0^\infty e^{-\lambda s}s^n\mathfrak Q(s)\,ds,
 \qquad n=0,1,\ldots .
 \tag{10.25}
\]
\end{hbtheorem}

\begin{hbtheorem}{thm:p15-w02b-p15-012-abstract-inversion}{proved=annex:p15-w02b-p15-012-abstract-inversion}{The abstract inversion formula (XIII.10, Theorem~3)}
For fixed \(t>0\), as \(n\to\infty\),
\[
 \frac{(-1)^{n-1}}{(n-1)!}\,\mathfrak R^{(n-1)}(n/t)\,(n/t)^n\longrightarrow\mathfrak Q(t).
 \tag{10.24}
\]
The left side is \(\int_0^\infty\mathfrak Q(s)\,G_{n,t}\{ds\}\) for the gamma
law \(G_{n,t}\) with density
\(e^{-ns/t}(ns/t)^{n-1}n/(t\,(n-1)!)\), which has expectation \(t\) and
variance \(t^2/n\). If \(\|\mathfrak Q(s)w-\mathfrak Q(t)w\|\leq\epsilon\)
whenever \(|s-t|<\delta\), the error at \(w\) is at most
\(\epsilon+2\|w\|t^2/(n\delta^2)\); the convergence is uniform for \(t\) in
every interval \(0<t\leq b\).
\end{hbtheorem}

Feller writes the limit in (10.24) as \(\lambda\to\infty\); the variable that
tends to infinity is \(n\). His proof is the one given in the annex: the left
side integrates the semigroup against a gamma law concentrating at \(t\)
\cite[vol.~II, XIII.10, pp.~462--463]{Feller1971ProbabilityII}. Applied to the
translation semigroup at a point \(x\), (10.24) is (6.6) for the function
\(s\mapsto w(x+s)\), now with convergence uniform in \(x\).

\begin{hbtheorem}{thm:p15-w02b-p15-012-resolvent-examples}{proved=annex:p15-w02b-p15-012-resolvent-examples}{The pseudo-Poisson resolvent and the translation boundary}
\begin{enumerate}
\item For the pseudo-Poisson semigroup (9.1) and \(\lambda>0\),
      \[
       \mathfrak R(\lambda)=\sum_{n=0}^\infty\frac{\alpha^n}{(\lambda+\alpha)^{n+1}}T^n,
       \qquad
       \|\lambda\mathfrak R(\lambda)-\mathbf 1\|\leq\frac{2\alpha}{\lambda+\alpha},
      \]
      so in this case (9.13) holds in operator norm.
\item For the translation semigroup on \(\mathscr L_u\),
      \(\mathfrak R(\lambda)w(x)=\int_0^\infty e^{-\lambda s}w(x+s)\,ds\), and
      for every \(t>0\) and \(n\geq1\) the operator
      \(E_{n,t}=\int_0^\infty T(s)\,G_{n,t}\{ds\}\) of (10.24) satisfies
      \(\|E_{n,t}-T(t)\|=2\). Thus (10.24) is strong convergence that is not
      convergence in operator norm.
\end{enumerate}
\end{hbtheorem}

\begin{hbwarning}{def:p15-w02b-p15-012-limits}{What the continuations do not give}
The generalized Taylor formula requires continuity and boundedness of \(u\) and
gives convergence for each fixed \(u\); \cref{thm:p15-w02b-p15-012-exponential-formula}
shows that it is not uniform over bounded functions. Formula (6.4) recovers
\(F\) only at continuity points; the remaining values follow from right
continuity. Formula (6.6) is uniform on intervals \(0<\theta\leq b\) and is
not asserted at \(\theta=0\), where \(n/\theta\) is undefined.
Both moment criteria are sufficient conditions, and no converse is asserted.
\Cref{thm:p15-w02b-p15-012-xv4-bridge} proves that (4.15) implies Carleman's
condition (\cref{thm:p15-criteria-carleman}); the reverse implication is not
proved, and no law separating the two conditions is constructed here.
The continuity theorems describe convergence only on bounded intervals of
continuity: mass may escape to infinity, the limit may be defective, and for
measures of infinite mass the transforms may diverge although the measures
converge (\cref{thm:p15-w02b-p15-012-laplace-continuity-examples}). For
characteristic functions, continuity of the limit at the origin cannot be
dropped (\cref{thm:p15-w02b-p15-012-cf-examples}). The inversion formulas
(3.5) and (3.11) need an integrability hypothesis on the transform; the
smoothed form (3.12) needs none but recovers increments of \(F\) only between
points of continuity. The values of a Laplace transform at the positive
integers determine the law (\cref{thm:p15-w02b-p15-012-laplace-integer-values}),
but nothing is asserted here about other sequences of points. P\'olya's
criterion is sufficient only, and its hypotheses cannot be weakened to those
printed in Feller's example (\cref{thm:p15-w02b-p15-012-polya-feller}). The exponential formula
(9.4) and the inversion (10.24) assert strong convergence only; for the
translation semigroup neither holds in operator norm
(\cref{thm:p15-w02b-p15-012-semigroup-examples,thm:p15-w02b-p15-012-resolvent-examples}).
Both require a continuous semigroup of contractions.
\end{hbwarning}
